%!TEX encoding = UTF-8 Unicode
%!TEX root = ../lect-w01-dod.tex

%%%%%%%%%%%%%%%%%%%%%%%%%%%%%%%%%%%%%%

%%%BEGIN TODO: flytta dessa till alla cls-filer compendium.cls lecturenotes.cls etc
% \newcommand{\ii}{\item}
% \newcommand{\is}{\item}
% \newcommand{\di}{\item}
% \newcommand{\ts}{\textit}
% \newcommand{\enquote}[1]{''#1''}
% \newcommand{\halfblankline}{\vspace{0.5em}}
% \newcommand{\blankline}{\vspace{1.0em}}
% \newcommand{\commandchar}{\texttt}
%%%END TODO
%%%%copypejsted from lect-preamble-dod 


\Subsection{Om Datorer och datoranvändning}

\begin{Slide}{Välkommen till Datorer och datoranvändning (dod)}
	\begin{itemize}
		\item Syftet med \emph{dod}-delen?
		      \begin{itemize}
			      \item Lär dig grundläggande verktyg.
			      \item Introduktion inför kommande kurser.
		      \end{itemize}
		\item Vem läser kursen?
		      \begin{itemize}
			      \item C-programmet.
			      \item D-programmet.
		      \end{itemize}
		\item Vad innehåller kursen?
		      \begin{itemize}
			      \item Fyra ämnen, med vardera en föreläsning och en laboration
			            \begin{enumerate}
				            \item Unix och Linux
				            \item \LaTeX
				            \item Git och Github
				            \item Maskinkod
			            \end{enumerate}
		      \end{itemize}
	\end{itemize}
\end{Slide}


\ifkompendium\else  %%%%%%%%%%%%%%%%%%%%%%%%%%%%%%%%%%


\begin{Slide}{Schema och gruppindelning}

	\begin{itemize}
		\item Gruppindelning
		      \begin{itemize}
			      \item Gruppindelning, se Canvas.
			      \item Kapten Alloc \url{https://fileadmin.cs.lth.se/pgk/kaptenalloc/}
		      \end{itemize}
		\item Schema
		\item Föreläsning alltid samma tid (tisdag 15:15), E:A
		\item Laborationer
		      \begin{itemize}
			      \item En labb per vecka, start på torsdag!
			      \item Föreberedelseuppgifter, kontrollfrågor.
			      \item Laboration 1, Linux (Första laborationen på torsdag!)
			      \item Laboration 2. \LaTeX
			      \item Laboration 3, Git
			      \item Laboration 4, Maskinkod
		      \end{itemize}
	\end{itemize}
\end{Slide}

\begin{Slide}{Kurshemsidan och Github}

	{\bf Den öppna kurshemsidan:}

	\smallskip
	\small\url{https://lunduniversity.github.io/pgk/}


	\bigskip


	{\bf Allt kursmaterial finns tillgängligt på GitHub:}


	\smallskip
	\begin{itemize}
		\item Allt material är \Emph{öppen källkod}.
		\item Skapa gärna ärenden \Eng{issues} om du upptäcker problem eller vill föreslå förbättringar.
		\item Förfrågningar om införande av implementerade ändringsförslag \Eng{pull requests} är hjärtligt välkomna! Öppna gärna först ett ärende för diskussion av dina ändringsförslag.
		\item Läs mer i kompendiets inledande kapitel om hur du bidrar till kursmaterialet.
	\end{itemize}
\end{Slide}

\begin{Slide}{Vid sjukdom / förhinder}

	Om du blir sjuka eller får annat förhinder:
	\begin{itemize}
		\item Meddela mig (dod) \url{mattias.nordahl@cs.lth.se}
		\item Meddela Björn (prog) \url{bjorn.regnell@cs.lth.se}
		\item Få hjälp och redovisa på resurstid
	\end{itemize}
\end{Slide}
\begin{Slide}{Vid frågor}

	\begin{itemize}
		\item Ni kan alltid maila mig: \url{mattias.nordahl@cs.lth.se}
		\item Discord. Hitta invite-länk på Canvas-sidan.
	\end{itemize}
\end{Slide}


\fi  %%%%%%%%%%%%%%%%%%%%%%%%%%%%%%%%%%%%%%%%%%%%%%%%%%%

\Subsection{Operativsystem}

\begin{Slide}{Förberedelse Lab 1: \texttt{\LabWeekONEdod}}
    \label{unix}
        \begin{itemize}
            \ii{Operativsystem, Unix historia, Linux}
            \is{Filer och kataloger}
            \is{Kommandon}
            \is{Filskydd}
            \is{Kommandotolk}
            \is{Processer}
        \end{itemize}
\end{Slide}

\begin{Slide}{Operativsystem}
    \begin{itemize}
        \ii{Dator --- kan köra program}
        \ii{Program --- instruktioner för datorn}
        \ii{Operativsystem --- en samling program som gör det möjligt att köra \enquote{vanliga} program}
        \ii{Operativsystemet hanterar:}
        \begin{itemize}
            \is{de program som körs (program körs ofta parallellt, operativsystemet ser till att programmen får minne och tid att exekvera)}
            \ii{yttre enheter (tangentbord, mus, skärm, nätverk, \ldots)}
            \ii{lagring av data (filer, \ldots)}
            \ii{skydd, felhantering, kommunikation med användaren}
        \end{itemize}

        \ii{Exempel:}
        \begin{itemize}
            \is{Linux, Windows, Mac OS X, Unix, \ldots}
        \end{itemize}
    \end{itemize}
\end{Slide}
\begin{Slide}{Unix historia}
    \begin{itemize}
        \ii{1960-talet: Multics, ett stort projekt avsedd för stora datorer.}
        \ii{1969: Ken Thompson på AT\&T Bell Labs utvecklar Unix, ett enklare system för mindre datorer med inspiration från Multics.}
        \ii{1973: Unix omskrivs i C av Dennis Ritchie, vilket underlättar distribution och anpassning.}
        \ii{Utvecklingen fortsätter med många varianter: Linux, BSD, Mac OS X, och Android bland andra.}
        \ii{Unix varumärket ägs av The Open Group som certifierar Unix-kompatibla system (Linux är ej certifierat).}
    \end{itemize}
\end{Slide}

\begin{Slide}{Linux och GNU}
    \begin{itemize}
        \ii{1983: GNU-verktygen -- fritt Unix-liknande operativsystem.}
        \ii{1991: Linus Torvalds skapar Linuxkärnan -- hanterar hårdvaruinteraktioner.}
        \ii{Linuxkärnan + GNU-verktyg = GNU/Linux -- skapar en fullständig användarmiljö.}
        \ii{\enquote{Linux} används ofta som kortform för hela systemet.}
    \end{itemize}
\end{Slide}

\begin{Slide}{Kännetecken för Unix}
    \textbf{Grundläggande:}
    \begin{itemize}
        \is{Kärna --- liten, grundläggande funktioner}
        \ii{Processer, tidsdelning}
        \ii{Filskydd}
        \ii{Fleranvändarsystem}
        \ii{Kommandotolk (terminalfönster)}
    \end{itemize}
    \textbf{Unix är:}
    \begin{itemize}
        \is{Litet, standardiserat, flyttbart}
        \ii{(Ursprungligen) Inte för nybörjare. Enklare nu med grafiska skrivbordsmiljöer.}
    \end{itemize}
\end{Slide}

\ifkompendium\else
\begin{Slide}{Undersökning OS}
    \textbf{Vilket operativsystem använder du?}

    \textit{Mentimeter eller handuppräckning}
\end{Slide}

\fi 

\begin{Slide}{Filer i Unix/Linux}

\begin{itemize}
\item  I en dator måste man kunna lagra
\begin{itemize}\SlideFontSmall
    \ii{program (Word, Excel, Java-/Scalakompilator, spelprogram, \ldots)}
    \ii{data (foton, document, programmeringsuppgifter, \ldots)}
\end{itemize}

\item Program och data lagras i \emph{filer}, som har ett namn:
\begin{itemize}\SlideFontSmall
  \item Filer ligger alltid i en \Emph{katalog} \Eng{directory}, även kallad \Emph{mapp} \Eng{folder}.
  \item Filer och kataloger lagras på \enquote{permanent} minne (hårddisk, SSD, USB-minne, \ldots).
\end{itemize}

\item Filnamn i Unix/Linux:
\begin{itemize}\SlideFontSmall
    \is{\texttt{namn.tillägg}\\Exempel: \texttt{Calc.scala}, \texttt{Calc.class}, \texttt{brev.txt}, \texttt{README.md},  \ldots}
    \ii{Tillägget är bara en extra upplysning; filtypen bestäms av det program som läser filen.}
    \ii{Filnamn kan sakna tillägg, exempel: \texttt{test1}}
\end{itemize}

\end{itemize}

\end{Slide}

\begin{Slide}{Kataloger}
    Varje fil finns i en katalog. Kataloger kan innehålla underkataloger, så man får en hierarkisk struktur, ett \Emph{filträd}:
    \begin{center}
        \resizebox{!}{0.5\textwidth}{%!TEX encoding = UTF-8 Unicode
%!TEX root = ../lect-w01-dod.tex

%%%%%%%%%%%%%%%%%%%%%%%%%%%%%%%%%%%%%%

% \documentclass{standalone}
% \usepackage{forest}

% \begin{document}

\begin{forest}
    for tree={
    font=\ttfamily,
    grow=south, % Tree grows down
    s sep=8mm, % Horizontal space between siblings
    l sep=6mm, % Vertical space between parent and child
    child anchor=north,
    parent anchor=south,
    anchor=north,
    align=center,
    % edge path={
    %   \noexpand\path [draw, \forestoption{edge}]
    %   (!u.parent anchor) -- ([yshift=2mm].child anchor) -| (.child anchor) \forestoption{edge label};
    % },
    edge path={
            \noexpand\path [draw, \forestoption{edge}]
            (!u.parent anchor) |- ([yshift=2mm].child anchor) -| (.child anchor) \forestoption{edge label};
        },
    draw, % Draw the border of nodes
    rounded corners=3mm, % Round corners of the nodes
    fill=white, % Fill color of nodes
    % drop shadow, % Optional: adds a shadow to nodes
    minimum height=6mm, % Minimum height of nodes
    minimum width=10mm, % Minimum width of nodes
    inner xsep=2mm, % X Padding within nodes
    inner ysep=1mm, % Y Padding within nodes
    calign=center
    }
    [/
    [\(\ldots\),draw=none]
    [usr]
    [home
        [\(\ldots\),draw=none]
        [anna
          [\(\ldots\),draw=none]
          [Downloads]
          [pgk
              [lab1
                  [kojo.scala,draw=none]
              ]
              [lab2]
              [\(\ldots\),draw=none]
          ]
          [dod
              [lab1
                  [\(\ldots\),draw=none]
                  [example.txt,draw=none]
                  [\(\ldots\),draw=none]
              ]
              [lab2]
              [\(\ldots\),draw=none]
          ]
          [\(\ldots\),draw=none]
        ]
        [\(\ldots\),draw=none]
    ]
    [bin]
    [etc]
    [\(\ldots\),draw=none]
    ]
\end{forest}

% \end{document}
}
    \end{center}
\end{Slide}

% \begin{Slide}{Filnamn}
%     I praktiken är det inte så rörigt som det ser ut på den förra bilden!
%     Ett fullständigt (absolut) filnamn börjar från rotkatalogen och man räknar upp alla katalogerna med \texttt{/} emellan:

%     \begin{Code}[language=]
%         /h/d24/b/uw3307bq-s/dod/lab1/example.txt
%     \end{Code}

%     Men operativsystemet håller reda på en \emph{aktuell katalog}, som ungefär betyder \emph{\enquote{var i filträdet befinner jag mig just nu?}}. Istället för att ange det absoluta filnamnet, räcker det att ange vägen med start från den aktuella katalogen -- ett s.k. \emph{relativt} filnamn.

%     \halfblankline

%     Om \texttt{uw3307bq-s} (användarens \emph{hemkatalog}) är aktuell katalog så kommer man åt filen ovan med namnet:

%     \begin{Code}[language=]
%         dod/lab1/example.txt
%     \end{Code}

% \end{Slide}


\begin{Slide}{Filnamn och sökvägar}\SlideFontSmall

\textbf{Termer:}
\begin{itemize}
\is{Sökväg (eng. path) -- vägen till en fil}
\ii{Absolut sökväg -- från rotkatalogen}
\ii{Relativ sökväg -- från aktuell katalog}
\ii{Aktuell katalog -- \emph{där du är i filträdet}}
\end{itemize}

\blankline\noindent
\textbf{Absolut filnamn:}
\begin{Code}[language=]
/home/anna/dod/lab1/example.txt
\end{Code}

\noindent\textbf{Relativt filnamn:}
\begin{Code}[language=]
dod/lab1/example.txt
\end{Code}
{\SlideFontTiny
På skolans datorer ser de översta delarna av filträdet annorlunda ut; en absolut sökväg för en användare med student-id \texttt{an1337sv-s} kan se ut ungefär såhär: \\\code{/h/d5/r/an1337sv-s/dod/lab1/example.txt}
}
\end{Slide}

% \begin{frame}[fragile]
%     \frametitle{Förkortade filnamn}
%     En del filer har förkortade namn:

%     \blankline
%     \begin{tabular}{cl}
%         \url{~}     & (tilde) hemkatalogen för aktuell användare (Alt Gr + $\sim$ ) \\
%         \url{~user} & hemkatalogen för användaren med användarnamnet user           \\
%         \texttt{.}    & (punkt) aktuell katalog                                       \\
%         \texttt{..}   & (punktpunkt) katalogen ovanför aktuell katalog i filträdet    \\
%     \end{tabular}

%     \blankline
%     Antag att \texttt{\url{~}/dod/lab1} är aktuell katalog. Man kan navigera upp och ned i filträdet enligt följande exempel:

%     \begin{Code}[language=]{6}
%         ~/pgk/lab1/kojo.scala
%         (till hemkatalogen, ner i pfk, ner i lab1)

%         ../../pgk/lab1/kojo.scala
%         (upp ett steg, upp ett steg, ner i pgk, ner i lab1)
%     \end{Code}

% \end{Slide}
\begin{Slide}{Förkortade filnamn}\SlideFontSmall

    \textbf{Förkortningar:}

    \ts{Vissa filnamn har speciella betydelser:}

    \blankline
        \begin{tabular}{cl}
            \texttt{\textasciitilde}     & Hemkatalog (tilde)              \\
            \texttt{\textasciitilde{}user} & Användarens hemkatalog (tilde user)          \\
            \texttt{.}    & Aktuell katalog (punkt)         \\
            \texttt{..}   & Överordnad katalog (punktpunkt) \\
        \end{tabular}
    \blankline

    \noindent\textbf{Exempel:}

    \ts{Om \texttt{\textasciitilde{}/dod/lab1} är aktuell katalog:}

\begin{Code}[language=]
~/pgk/lab1/kojo.scala       
\end{Code}
Betydelse: till hemkatalogen, ner i pfk, ner i lab1
\begin{Code}[language=]
../../pgk/lab1/kojo.scala
\end{Code}
Betydelse: upp ett steg, upp ett steg, ner i pgk, ner i lab1

\end{Slide}

% \begin{frame}[fragile]
%     \frametitle{Kommandotolk}
%     Man kommunicerar med Unix eller Linux genom att ge kommandon. Ett särskilt program i operativsystemet, kommandotolken, ''shell'', läser in kommandona och utför det som ska göras.

%     \pindent Det finns flera olika kommandotolkar. På Linuxdatorerna i E-huset används kommandotolken \texttt{bash} (''Bourne-Again SHell''). Andra vanliga kommandotolkar är \texttt{sh}, \texttt{zsh} och \texttt{csh}.

%     \pindent Kommandotolken läser kommandon som ges i ett terminalfönster på skärmen. Så här arbetar tolken:

%     \begin{Code}[language=]{4}
%         Upprepa i all oändlighet:
%             Läs ett kommando
%             Om kommandot är ett riktigt kommando:
%                 utför det som ska göras
%             annars:
%                 skriv ut felmeddelande
%     \end{Code}

% \end{Slide}
\begin{Slide}{Kommandotolk (shell, ''skal'') och terminal}

    \textbf{Termerna \emph{kommandotolk} och \emph{terminal} används ofta synonymt, men har olika betydelser:}

    \halfblankline
    \begin{itemize}
        \item \textbf{Kommandotolk}: {För kommunikation med operativsystemet via kommandon}
        \item \textbf{Terminal}: {Program som låter användaren skriva kommandon till kommandotolken}
        \item \textbf{Shell}: {Engelska för kommandotolk}, \textit{skal} på svenska
    \end{itemize}

    \blankline

    \noindent\textbf{Exempel på kommandotolkar:}
    \begin{itemize}
        \is{\texttt{bash} (Bourne-Again SHell)}
        \is{\texttt{sh}, \texttt{zsh}, \texttt{csh}}
    \end{itemize}

\end{Slide}



\begin{Slide}{Kommandotolk (shell, ''skal'') och terminal}

Så här fungerar en kommandotolk:

\begin{Code}[language=]
Upprepa i all oändlighet:
  Läs ett kommando
  Om kommandot är ett riktigt kommando:
    utför det som ska göras
  annars:
    skriv ut felmeddelande
\end{Code}

\vfill\url{https://youtu.be/zQ-e9SuoWXo}

\end{Slide}


\begin{Slide}{Bash vs. Zsh för Mac-användare}

    \noindent\textbf{Bash och Zsh:}
    \begin{itemize}
        \is{Bash är standard i de flesta Linux-system.}
        \is{Zsh är standard i macOS från Catalina.}
        \is{Båda är kraftfulla skal i Unix-system.}
        \is{För grundläggande användning, t.ex. denna kurs, fungerar båda bra.}
    \end{itemize}

    \noindent\textbf{Att byta till Bash (valfritt):}
    \begin{itemize}
        \is{Öppna Terminal och skriv \texttt{bash} för att temporärt byta till Bash.}
        \is{Installera senaste Bash via Homebrew: \texttt{brew install bash}}
    \end{itemize}

    \noindent\textbf{Skillnader:}
    \begin{itemize}
        \is{Zsh upplevs av många som mer användarvänligt. Erbjuder fler funktioner och anpassningsmöjligheter, men kan vara överväldigande för nybörjare.}
        \is{Bash är mer konservativt och har bättre kompatibilitet.}
    \end{itemize}

\end{Slide}
\begin{Slide}{Unix-shellalternativ för Windows-användare}

    \noindent\textbf{Git Bash:}
    \begin{itemize}
        \is{Enkel Bash-tillgång, minimal konfiguration.}
        \is{Installeras med Git för Windows.}
    \end{itemize}

    \noindent\textbf{WSL (Windows Subsystem for Linux):}
    \begin{itemize}
        \is{Kör fullständig Linux i Windows.}
    \end{itemize}

    \noindent\textbf{Cygwin:}
    \begin{itemize}
        \is{Ger bred Unix-vertygsuppsättning.}
        \is{Mer konfiguration.}
    \end{itemize}

    \blankline
    \noindent\textbf{Vilket ska jag välja?}
    \begin{itemize}
        \is{\textbf{Git Bash}: Grundläggande Bash och Git. (Rekommenderas)}
        \is{\textbf{WSL}: Fullständig Linux-utveckling. (Rekommenderas)}
        \is{\textbf{Cygwin}: Omfattande Unix-funktionalitet. (Avråds)}
    \end{itemize}

\end{Slide}

\begin{Slide}{Kommandoformat}\SlideFontSmall
    \textbf{Varje kommando skrivs på en rad:}

\begin{Code}[language=]
kommandonamn -option1 -option2 ... argument1 ...
\end{Code}

    \begin{itemize}
        \ii{Kommandot talar om vad som ska göras.}
        \ii{Optionerna modifierar kommandot på något sätt.}
        \ii{Argumenten är (oftast) filer eller kataloger som påverkas av kommandot.}
    \end{itemize}

    \halfblankline
    \noindent\textbf{Exempel:}

    \halfblankline
        \begin{tabular}{lp{60mm}}
        \texttt{scala run Calc.scala}        & Kompilera och kör \texttt{Calc.scala}\\
        \texttt{scala run .}        & Kompilera och kör alla Scala-filer i aktuell katalog och alla dess underkataloger\\
        \texttt{cp -i report.tex old.tex} & Kopierar \texttt{report.tex} till \texttt{old.tex}. \texttt{-i} betyder att systemet frågar om \texttt{old.tex} ska skrivas över, om den redan finns. \\
        \end{tabular}
\end{Slide}



\begin{Slide}{Kommandon för filhantering}\SlideFontSmall
    Exempel på kommandon för att hantera filer:

    \blankline
    \begin{tabular}{lp{7.0cm}}
        \texttt{cp orig kopia} & Kopiera filen \texttt{orig} till \texttt{kopia}.                                                                                                                                 \\
        \texttt{less fil}      & Skriv ut \texttt{fil} på skärmen, en sida i taget.                                                                                                                             \\
        \texttt{ls [-la] kat}  & Skriv ut en innehållsförteckning över katalogen \texttt{kat} (aktuell katalog om ingen katalog ges). \texttt{-l} skriv i ett längre format, \texttt{-a} skriv också ut punktfiler. \\
        \texttt{mv fil1 fil2}  & Döp om \texttt{fil1} till \texttt{fil2}. Om \texttt{fil2} är en katalog så flyttas filen till den katalogen.                                                                       \\
        \texttt{rm fil1 ...}   & Tag bort de angivna filerna.                                                                                                                                                 \\
    \end{tabular}
\end{Slide}

\begin{Slide}{Kataloghantering}\SlideFontSmall
    Exempel på kommandon för att hantera kataloger:

    \blankline
    \begin{tabular}{lp{8.3cm}}
        \texttt{cd kat}    & Ändra aktuell katalog till \texttt{kat}. Utan argument blir det hemkatalogen (samma som \texttt{cd} \url{~}). \\
        \texttt{mkdir kat} & Skapa underkatalogen \texttt{kat} i den aktuella katalogen.                                                 \\
        \texttt{pwd}       & Skriv ut namnet på aktuell katalog.                                                                       \\
        \texttt{rmdir kat} & Tag bort \texttt{kat}. Man måste först ta bort alla filerna i katalogen.                                    \\
    \end{tabular}
\end{Slide}
% \begin{frame}[fragile]
%     \frametitle{Filskydd 1}
%     Varje fil har en ägare. Ägaren identifieras med användarnamn och grupp, till exempel \texttt{ma1337no-s} i gruppen \texttt{students}.

%     \pindent Ägaren kan skydda filer så att andra inte kan komma åt dem, eller göra dem publika så att andra kan komma åt dem.

%     \pindent För att ta reda på filskydd och ägare (och storlek och datum) för filer kan man använda kommandot \texttt{ls} med optionen \texttt{l}, alltså \texttt{ls -l}:
%     \begin{Code}[language=]{6}
%         hacke-3{ma1337no-s}: ls -l
%         -rw-r-----  1 ma1337no-s  students  4940 aug 21 11:14 example.txt
%          (skydd)      (ägare)   (grupp)
%     \end{Code}

% \end{Slide}

\begin{Slide}{Filskydd 1/3}

    \noindent\textbf{Filägande:}
    \begin{itemize}
        \is{Ägare identifieras via användarnamn och grupp\\
            \hspace{1em}(t.ex., \texttt{ma1337no-s} i \texttt{students})}
    \end{itemize}

    \noindent\textbf{Filsäkerhet:}
    \begin{itemize}
        \is{Ändra åtkomsträttigheter för att skydda eller dela filer}
    \end{itemize}

    \noindent\textbf{Visa detaljer:}
\begin{Code}[language=]
hacke-3{ma1337no-s}: ls -l
-rw-r-----  1 ma1337no-s  students  4940 aug 21 11:14 example.txt
(skydd)       (ägare)     (grupp)
\end{Code}

\end{Slide}

\begin{Slide}{Filskydd 2/3}
\begin{Code}[language=]
-rw-r-----  1 ma1337no-s  students  4940 aug 21 11:14 example.txt
u  g  o
\end{Code}

    Tre kategorier av användare:

    \begin{description}
        \item[\texttt{u}] (user) filens ägare
        \item[\texttt{g}] (group) medlemmar i samma grupp som ägaren
        \item[\texttt{o}] (others) alla andra
    \end{description}

    Tre olika rättigheter:

    \begin{description}
        \item[\texttt{r}] (read) tillstånd att läsa
        \item[\texttt{w}] (write) tillstånd att skriva
        \item[\texttt{x}] (execute) tillstånd att exekvera program (för kataloger betyder \texttt{x} tillstånd att titta på innehållet i katalogen)
    \end{description}

\end{Slide}
\begin{Slide}{Filskydd 3/3}
\textbf{Kommando för att ändra filskydd:}

\begin{Code}[language=]
    chmod skydd fil1 fil2 etc
\end{Code}

\noindent\textbf{Filskydd kan anges symboliskt, till exempel \texttt{u=rw,o=r}}

\halfblankline
\noindent\textbf{Eller numeriskt, enligt:}

\begin{Code}[language=]
    x = 001, w = 010, r = 100, rx = 101, rw = 110, rwx = 111
\end{Code}


\textbf{Siffrorna tolkar man sedan som binära tal, vilket ger:}

\begin{Code}[language=]
    x = 1, w = 2, r = 4, rx = 5, rw = 6, rwx = 7
\end{Code}


\noindent\textbf{Exempel numeriska filskydd och motsv. symboliska:}

\begin{Code}[language=]
 chmod 604 fil1 fil2         chmod u=rw,o=r fil1 fil2   
 chmod 755 fil1 fil2         chmod u=rwx,g=rx,o=rx fil1 fil2
\end{Code}
\end{Slide}
\begin{Slide}{bash-finesser 1} \SlideFontSmall
    Man kan editera den aktuella kommandoraden:

    \blankline
    \begin{tabular}{p{2.2cm}p{8cm}}
        \commandchar{$\leftarrow$\ $\rightarrow$} & Flytta markören på raden.                                                 \\
        \commandchar{Control-K}                   & Radera hela raden.                                                        \\
        \commandchar{Control-U}                   & Radera allt till vänster.                                                 \\
        \commandchar{Control-L}                   & Töm terminalfönstret.                                                     \\
        \commandchar{Tab}                         & Filnamnskomplettering (man behöver bara skriva början av ett långt namn). \\
    \end{tabular}

    \blankline
    \texttt{bash} håller reda på de senaste kommandona som utförts, och man kan få tillbaka gamla kommandon:

    \blankline
    \begin{tabular}{p{2.2cm}p{8cm}}
        \texttt{$\uparrow$} \texttt{$\downarrow$} & Bläddra bland tidigare kommandon.                  \\
        \texttt{!abc}                           & Gör om senaste kommando som inleds med \texttt{abc}. \\
        \texttt{!!}                             & Gör om senaste kommando.                           \\
    \end{tabular}
\end{Slide}

\begin{Slide}{bash-finesser 2}
    \textbf{När man refererar till filer kan man utnyttja jokertecken (wildcards):}

    \blankline
        \begin{tabular}{ll}
            \texttt{?} & motsvarar \emph{ett} godtyckligt tecken \\
            \texttt{*} & motsvarar 0--flera godtyckliga tecken   \\
        \end{tabular}

    \blankline
    \textbf{Antag att vi har följande filer i vår katalog:}
\begin{Code}[language=]
Test1.java  Test2.java  Test23.java  Final.java
Test1.class Test2.class Final.class  FinalReport.tex
Outline.tex
\end{Code}

    \textbf{Då kan vi till exempel skriva:}
\begin{Code}[language=]
javac Test?.java (javac Test1.java Test2.java)
rm *.class       (rm Test1.class Test2.class Final.class)
gedit F*.tex     (gedit FinalReport.tex)
\end{Code}

\end{Slide}
% \begin{frame}[fragile]
%     \frametitle{Initieringsfiler}
%     Många program läser en initieringsfil när de startas. Initieringsfilen innehåller inställningar för programmet. Namnen på initieringsfilerna inleds med punkt, så man ser normalt inte dessa filer när man gör \texttt{ls}.

%     \pindent \texttt{bash} läser initieringsfilen \texttt{.bash\_profile} när ett nytt terminalfönster skapas. Bra rader att ha i denna fil (gör att man alltid får frågor om man verkligen vill ta bort/skriva över filer):

%     \begin{Code}[language=]
%         alias rm='rm -i'
%         alias cp='cp -i'
%         alias mv='mv -i'
%     \end{Code}

%     Varning: kopiera inte andras initieringsfiler om du inte begriper vad som står i dem!
% \end{Slide}
\begin{Slide}{Initieringsfiler}

    \noindent\textbf{Om initieringsfiler:}
    \begin{itemize}
        \is{Används ofta för att konfigurera program.}
        \ii{Filnamn börjar med punkt och är dolda, s.k. \emph{punktfiler}.}
    \end{itemize}

    \noindent\textbf{\texttt{bash} initieringsfil:}
    \begin{itemize}
        \is{bash läser \texttt{.bash\_profile} (från hemkatalogen) när ett nytt terminalfönster öppnas.}
        \ii{Nyttiga alias för att förhindra oavsiktliga filoperationer:}
    \end{itemize}

\begin{Code}[language=]
alias rm='rm -i'
alias cp='cp -i'
alias mv='mv -i'
\end{Code}

    \textbf{Varning:}
    \begin{itemize}
        \is{Kopiera inte initieringsfiler utan att förstå innehållet!}
    \end{itemize}

\end{Slide}

% \begin{frame}[fragile]
%     \frametitle{Vad är egentligen ett kommando?}
%     Kommandon är av två slag:

%     \begin{description}
%         \item[''inbyggda''] små kommandon som exekveras direkt i kommandotolken, till exempel \texttt{cd}, \texttt{pwd}, \texttt{exit}
%         \item[''vanliga program''] kommandotolken letar efter ett program med samma namn som kommandot
%     \end{description}

%     Kommandotolken letar efter program i de kataloger som anges i sökvägen (operativsystemvariabeln \texttt{PATH}, som man själv kan sätta med kommandot \texttt{export}). När man till exempel skriver \texttt{javac X.java} exekveras programmet \texttt{/usr/bin/javac}, eftersom katalogen \texttt{/usr/bin} finns i sökvägen.

%     \pindent När man ska exekvera ett eget program som finns i den aktuella katalogen skriver man:
%     \begin{Code}[language=]
%         ./programnamn
%     \end{Code}

%     \texttt{./} är nödvändigt eftersom \texttt{.} (den aktuella katalogen) inte finns i sökvägen.
% \end{Slide}
\begin{Slide}{Vad är egentligen ett kommando?}

    \noindent\textbf{Typer av kommandon:}
    \begin{description}
        \di{Inbyggda}{Kör direkt i kommandotolken (t.ex. \texttt{cd}, \texttt{pwd}, \texttt{ls}).}
        \di{Vanliga program}{Kommandotolken söker efter programfiler i \texttt{PATH}.}
    \end{description}

    \noindent\textbf{Programexekvering:}
    \begin{itemize}
        \is{\texttt{PATH} definierar var program letas upp (\texttt{export PATH=...} för anpassning).}
        \ii{Ex: \texttt{/usr/bin/javac} finns i \texttt{PATH}, så \texttt{javac X.java} går att köra.}
    \end{itemize}

    \noindent\textbf{Köra egna program:}
\begin{Code}[language=]
./programnamn
\end{Code}
{\SlideFontSmall Använd \texttt{./} för program i den aktuella katalogen, eftersom \texttt{.} inte ingår i \texttt{PATH}.}

\end{Slide}

% \begin{frame}[fragile]
%     \frametitle{Processer}
%     Varje program som körs i Unix körs som en egen process som exekverar självständigt, ''parallellt'' med andra processer. Varje process körs några millisekunder, sedan får nästa process exekvera, osv. Detta hanteras av operativsystemet.

%     \pindent När man exekverar ett program från ett terminalfönster ''låser'' programmet fönstret tills det avslutas --- så vill man ofta att det ska fungera.

%     \pindent Men man kan också köra program ''i bakgrunden'' dvs frikopplade från terminalfönstret. Det gör man genom att skriva \texttt{\&} sist på kommandoraden, till exempel:

%     \begin{Code}[language=]
%         gedit example.txt &
%     \end{Code}

% \end{Slide}

\begin{Slide}{Processer}\SlideFontSmall

    \noindent\textbf{Grundläggande om Processer:}
    \begin{itemize}
        \ii{Varje program körs som en egen process, oberoende och ''parallellt'' med andra.}
        \ii{Processväxling sker kontinuerligt, hanteras automatiskt av operativsystemet.}
    \end{itemize}

    \noindent\textbf{Exekvering i Terminal:}
    \begin{itemize}
        \ii{Standard: Kommandotolken väntar på programmet. Terminalen framstår som upptagen tills programmet avslutas.}
        \ii{Bakgrundskörning: Använd \texttt{\&} för att köra programmet frikopplat från terminalen.}
    \end{itemize}

    \noindent\textbf{Exempel:}
\begin{Code}[language=]
gedit example.txt &
\end{Code}
    \textbf{\textit{Öppnar \texttt{gedit} i bakgrunden, terminalen blir omedelbart tillgänglig för nya kommandon.}}

\end{Slide}

% \begin{frame}[fragile]
%     \frametitle{Fönstersystem, skrivbordsmiljö}
%     I Unix och Linux är fönstersystemet inte inbyggt (Unix utvecklades långt innan det fanns grafiska skärmar).

%     \pindent De flesta Unix- och Linuxsystem som kör på persondatorer använder fönstersystemet X Window System, som oftast kallas bara X. X gör det möjligt att ha fönster på skärmen, att flytta dem, ikonifiera dem, osv.

%     \pindent Exakt hur fönstren ska se ut och hur man arbetar med dem bestäms av \emph{fönsterhanteraren}, som också är ett separat program. Ovanpå fönsterhanteraren finns skrivbordsmiljön.

%     \pindent Det bästa (enda) sättet att lära sig att använda en skrivbordsmiljö är att träna!
% \end{Slide}


\begin{Slide}{Fönstersystem, skrivbordsmiljö}\SlideFontSmall
    \noindent\textbf{Fönstersystem:}
    \begin{itemize}
        \ii{I Unix och Linux är fönstersystemet inte inbyggt (Unix utvecklades långt innan det fanns grafiska skärmar).}
        \ii{Ett äldre fönstersystem är X Window System (ofta kallat \emph{X}).}
        \ii{Ett nyare fönstersystem är Wayland (används i Ubuntu).}
    \end{itemize}

    \noindent\textbf{Fönsterhanterare och skrivbordsmiljö:}
    \begin{itemize}
        \ii{Fönsterhanteraren bestämmer utseende och funktion hos fönstren.}
        \ii{Skrivbordsmiljön bygger vidare på fönsterhanteraren för en fullständig användarupplevelse.}
    \end{itemize}

    \noindent\textbf{Lära sig skrivbordsmiljön:}
    \begin{itemize}
        \ii{För många av er är Linux nytt. Bästa sättet att bli bekant med en skrivbordsmiljö är att prova sig fram och träna!}
        \ii{I Ubuntu: Tryck på windowsknappen och skriv ''help'' så får du fram \Emph{Ubuntu Desktop Guide} med en översikt av skrivbordsmiljön.}
        \ii{Lär dig gärna kortkommandon för funktioner du använder extra ofta, se ''Tips \& Tricks -> Useful keyboard shortcuts'' i desktop-gajden.}
    \end{itemize}

\end{Slide}